% GENERATED FILE. Edit canonical lessons, then run npm run generate:books.
% canonical-source-hash: fnv1a64:6079d14592306807
% canonical-lessons: FA-C03-shoma-to, FA-C03-chist, FA-C03-esm-e-shoma-chist, FA-C03-khoshvaghtam, FA-C03-practice

\chapter{Ask and Answer Names}
\label{ch:fa-ask-and-answer-names}

This chapter is generated from the canonical micro-lessons used by Language Ladder.
Each section stays independently resumable and preserves the authored prerequisite order.

\section[shomâ / to]{\fa{شما} / \fa{تو} — two relationships inside “you”}

\label{lesson:FA-C03-shoma-to}



\begin{quote}
\textbf{Warm-up.} {[}PAUSE 2s{]} Say \textbf{esm-e man ... ast} with your own name. To turn that sentence around, you next need a safe word for the other person.
\end{quote}

\subsection*{The two words}
\noindent
\begin{tabularx}{\linewidth}{@{}>{\raggedright\arraybackslash}X>{\raggedright\arraybackslash}X>{\raggedright\arraybackslash}X@{}}
\toprule
\textbf{Persian} & \textbf{read} & \textbf{use} \\
\midrule
\textbf{\fa{شما}} & \emph{shomâ} & one person respectfully, or more than one person \\
\textbf{\fa{تو}} & \emph{to} & one familiar person \\
\bottomrule
\end{tabularx}

In \textbf{\fa{شما}}, meet \textbf{\fa{ش}} \emph{sh}; its three dots distinguish it from \textbf{\fa{س}} \emph{s}. Final \textbf{\fa{ا}} carries long \emph{â}. In \textbf{\fa{تو}}, \textbf{\fa{و}} has its \emph{o} value rather than the long \emph{u} it had in \textbf{\fa{ممنون}}. Learn the sound from the whole word.

\begin{grammarlens}[title={Grammar Lens: begin respectfully}]
Persian uses plural \textbf{shomâ} for respectful singular address, much as Russian uses \emph{vy} and French uses \emph{vous}. Use \textbf{shomâ} with an unfamiliar adult or in a first meeting. Use \textbf{to} after the relationship is familiar. English once made a similar distinction but now uses \emph{you} for both.
\end{grammarlens}

\begin{cousinweb}
Persian \textbf{to} belongs to the ancient Indo-European familiar-“you” family: English \textbf{thou}, Latin \emph{tū}, and Sanskrit \emph{tvam} are relatives. The word is old; the social choice carried by it is current.
\end{cousinweb}

\subsection*{Guided Practice}
\begin{itemize}
\raggedright
  \item {[}YOU SAY: a new adult — \textbf{shomâ}{]}
  \item {[}YOU SAY: a close friend — \textbf{to}{]}
  \item {[}YOU SAY: the old family — \textbf{to, thou, tū, tvam}{]}
\end{itemize}

\begin{tcolorbox}[breakable,colback=teal!4,colframe=teal!35!black,title={Wrap-up recall}]
Which form is safest when the relationship is new? (\textbf{Shomâ}.) Which is the relative of English \emph{thou}? (\textbf{To}.)

Source: Persian Online: Pronouns.
\end{tcolorbox}
\section[chist]{\fa{چیست} — “what is?” in one joined word}

\label{lesson:FA-C03-chist}



\begin{quote}
\textbf{Warm-up.} {[}PAUSE 2s{]} Which “you” is safe for a new adult? (\textbf{\fa{شما}}, \emph{shomâ}.) Now add the question engine that will follow “your name.”
\end{quote}

\subsection*{You'll want to know first — \fa{چیست}}
\begin{quote}
\textbf{\fa{چیست}} — \emph{chist} — \textbf{what is?}
\end{quote}

Read from the right: \textbf{\fa{چ}} \emph{ch}, \textbf{\fa{ی}} long \emph{i}, \textbf{\fa{س}} \emph{s}, \textbf{\fa{ت}} \emph{t}. \textbf{\fa{چ}} is one of Persian's four additions to the Arabic alphabet: it uses the same base shape as \textbf{\fa{ج}}, with three dots below. Meet only this joined form now.

\begin{grammarlens}[title={Grammar Lens: two words compressed}]
\textbf{chist} is the careful fusion of \textbf{chi} “what” and \textbf{ast} “is.” The opening vowel of \emph{ast} disappears when the pieces join:

\begin{quote}
\emph{chi + ast} $\to$ \textbf{chist}
\end{quote}

You already know full \textbf{ast} from \emph{esm-e man ... ast}. This lesson changes no copula rule; it gives one high-frequency fused question form.
\end{grammarlens}

\subsection*{Guided Practice}
\begin{itemize}
\raggedright
  \item {[}YOU SAY: \textbf{chi + ast $\to$ chist}{]}
  \item {[}YOU READ: \textbf{\fa{چیست}} from \textbf{\fa{چ}} to \textbf{\fa{ت}}{]}
  \item {[}YOU SAY: \textbf{chist?} — “what is?”{]}
\end{itemize}

\begin{tcolorbox}[breakable,colback=teal!4,colframe=teal!35!black,title={Wrap-up recall}]
What two pieces make \textbf{chist}? (\textbf{Chi + ast}.) Which new Persian letter opens it? (\textbf{\fa{چ}}, \emph{ch}, with three dots below.)
\end{tcolorbox}
\section[esm-e shomâ chist?]{\fa{اسم} \fa{شما} \fa{چیست؟} — ask a name respectfully}

\label{lesson:FA-C03-esm-e-shoma-chist}



\begin{quote}
\textbf{Warm-up.} {[}PAUSE 2s{]} Give the three ready pieces: \textbf{esm-e} “name-of,” \textbf{shomâ} respectful “you,” and \textbf{chist} “what is?”
\end{quote}

\subsection*{The exchange}
\begin{quote}
\textbf{\fa{اسم} \fa{شما} \fa{چیست؟}} — \emph{esm-e shomâ chist?} — \textbf{What is your name?}
\end{quote}

The Persian order is transparent once the pieces are known:

\begin{quote}
name + \textbf{-e} + you + what-is?
\end{quote}

The spoken ezafe turns \textbf{esm-e shomâ} into “your name.” \textbf{Chist} stays at the end. No new agreement or verb table is hiding here.

\subsection*{Script: the question mark}
Persian ends a written question with \textbf{\fa{؟}}. Its curve faces the opposite way from English \textbf{?}, matching the right-to-left line. Read the sentence from its right edge; meet the punctuation at the far left.

\subsection*{Guided Practice}
\begin{itemize}
\raggedright
  \item {[}YOU SAY: \textbf{esm-e shomâ chist?}{]}
  \item {[}YOU SAY: ask, then answer — \textbf{esm-e man ... ast}{]}
  \item {[}YOU POINT: to \textbf{\fa{؟}} at the end of the right-to-left question{]}
\end{itemize}

\begin{tcolorbox}[breakable,colback=teal!4,colframe=teal!35!black,title={Wrap-up recall}]
Ask the question once without reading the romanization. Then answer with the careful name frame you already know.

Source: Persian Online: Introducing Yourself.
\end{tcolorbox}
\section[khoshvaghtam]{\fa{خوشوقتم} — pleased to meet you}

\label{lesson:FA-C03-khoshvaghtam}



\begin{quote}
\textbf{Warm-up.} {[}PAUSE 2s{]} Ask \textbf{esm-e shomâ chist?} and answer \textbf{esm-e man ... ast}. One warm response now closes that tiny meeting.
\end{quote}

\subsection*{The exchange}
\begin{quote}
\textbf{\fa{خوشوقتم}} — \emph{khoshvaghtam} — \textbf{Pleased to meet you.}
\end{quote}

Treat the joined word as one social move first. \textbf{\fa{خ}} gives \emph{kh}, like the sound in Scottish \emph{loch}. \textbf{\fa{ش}} gives \emph{sh}. Persian commonly pronounces \textbf{\fa{ق}} here with the merged \emph{gh/q} sound of contemporary Iranian speech.

\begin{cousinweb}
\begin{itemize}
\raggedright
  \item \textbf{khosh} — pleasant, happy; inherited Persian vocabulary.
  \item \textbf{vaqt} (pronounced \emph{vaght} here) — time; an Arabic loan fully at home in Persian.
  \item \textbf{-am} — I am, the first-person ending.
\end{itemize}

So the word builds “pleasant-time-I-am” into the idiomatic feeling “I am pleased.” The history is layered; the expression is ordinary modern Persian.
\end{cousinweb}

\subsection*{Guided Practice}
\begin{itemize}
\raggedright
  \item {[}YOU SAY: \textbf{khoshvaghtam} as one smooth response{]}
  \item {[}YOU SAY: the three pieces — \textbf{khosh + vaqt + -am}{]}
  \item {[}YOU SAY: which piece came from Arabic — \textbf{vaqt}{]}
\end{itemize}

\begin{tcolorbox}[breakable,colback=teal!4,colframe=teal!35!black,title={Wrap-up recall}]
Which response closes the introduction? (\textbf{Khoshvaghtam}.) Which three layers can you hear inside it? (\textbf{Khosh + vaqt + -am}.)

Source: Persian Online: Introducing Yourself.
\end{tcolorbox}
\section[Practice]{Practice — a complete Persian name exchange}

\label{lesson:FA-C03-practice}



\begin{quote}
\textbf{Warm-up.} {[}PAUSE 2s{]} Recall the four moves without adding anything: greet, ask, answer, acknowledge.
\end{quote}

\subsection*{The exchange}
\begin{quote}
A: \textbf{\fa{سلام}!} — \emph{Salâm!} B: \textbf{\fa{سلام}!} — \emph{Salâm!} A: \textbf{\fa{اسم} \fa{شما} \fa{چیست؟}} — \emph{Esm-e shomâ chist?} B: \textbf{\fa{اسم} \fa{من} \fa{سارا} \fa{است}.} — \emph{Esm-e man Sârâ ast.} A: \textbf{\fa{خوشوقتم}.} — \emph{Khoshvaghtam.} B: \textbf{\fa{ممنون}.} — \emph{Mamnun.}
\end{quote}

Every line is already known. \textbf{Shomâ} keeps the first meeting respectful; the answer returns to \textbf{man}. The question mark sits at the left end of the right-to-left question.

\subsection*{Guided Practice}
\begin{itemize}
\raggedright
  \item {[}YOU SAY: both voices, using your own name{]}
  \item {[}YOU SAY: the exchange again without romanization{]}
  \item {[}YOU SAY: only the three name moves — question, answer, response{]}
\end{itemize}

\begin{tcolorbox}[breakable,colback=teal!4,colframe=teal!35!black,title={Wrap-up recall}]
Run the meeting once from \textbf{salâm} to \textbf{mamnun}. If one line stalls, review that single micro-lesson rather than rereading the whole chapter.
\end{tcolorbox}
